% ===== PRÉAMBULE CANONIQUE — à coller VERBATIM en tête de CHAQUE .tex =====
\documentclass[11pt,a4paper]{article}
\usepackage[utf8]{inputenc}\usepackage[T1]{fontenc}\usepackage{lmodern}
\usepackage[french]{babel}
\usepackage{amsmath,amssymb,mathtools}\usepackage{icomma}
\usepackage{siunitx}\sisetup{locale=FR}  % \qty{}{}, \si{}, \ang{}
\usepackage{xcolor}\usepackage{colortbl}
\usepackage{tikz}\usepackage{pgfplots}\pgfplotsset{compat=1.18}
\usepackage[siunitx,europeanresistors]{circuitikz} % circuits (élec.)
\usepackage[most]{tcolorbox}\usepackage{enumitem}\usepackage{array,booktabs}
\usepackage[a4paper,margin=2.2cm]{geometry}
\usetikzlibrary{arrows.meta,calc,angles,quotes,positioning,patterns,%
  backgrounds,shapes.geometric,decorations.pathmorphing,decorations.markings,intersections}
\definecolor{primary}{RGB}{222,49,99}\definecolor{primarydark}{RGB}{150,22,55}
\definecolor{defcol}{RGB}{0,114,178}\definecolor{methcol}{RGB}{52,92,148}
\definecolor{excol}{RGB}{0,135,90}\definecolor{attncol}{RGB}{200,40,55}
\tcbset{boxbase/.style={enhanced,breakable,boxrule=0.9pt,arc=2.5pt,
  left=3mm,right=3mm,top=2.3mm,bottom=2.3mm,fonttitle=\bfseries,coltitle=white}}
% --- boîtes TUTORIEL (noms EXACTS, ne pas renommer) ---
\newtcolorbox{cle}[1][]{boxbase,colback=defcol!6,colframe=defcol,
  title={\textsc{À retenir}\ifx\relax#1\relax\relax\else\textmd{ : \itshape #1}\fi}}
\newtcolorbox{methode}[1][]{boxbase,colback=methcol!7,colframe=methcol,sharp corners,
  title={\textsc{Méthode}\ifx\relax#1\relax\relax\else\textmd{ --- \itshape #1}\fi}}
\newtcolorbox{exemplebox}[1][]{boxbase,colback=excol!5,colframe=excol,
  title={\textsc{Exemple}\ifx\relax#1\relax\relax\else\textmd{ : \itshape #1}\fi}}
\newtcolorbox{piege}[1][]{boxbase,colback=attncol!6,colframe=attncol,
  title={\textsc{Piège}\ifx\relax#1\relax\relax\else\textmd{ : \itshape #1}\fi}}
\newtcolorbox{remarque}[1][]{boxbase,colback=black!4,colframe=black!45,
  title={\textsc{Remarque}\ifx\relax#1\relax\relax\else\textmd{ : \itshape #1}\fi}}
% --- boîtes EXERCICE / CORRIGÉ (noms EXACTS, auto-comptées) ---
\newtcolorbox[auto counter]{exo}[1][]{boxbase,colback=primary!4,colframe=primary,
  title={\textsc{Exercice}~\thetcbcounter\ifx\relax#1\relax\relax\else\textmd{ --- \itshape #1}\fi}}
\newtcolorbox[auto counter]{corrige}[1][]{boxbase,colback=excol!4,colframe=excol,
  title={\textsc{Corrigé}~\thetcbcounter\ifx\relax#1\relax\relax\else\textmd{ --- \itshape #1}\fi}}
\newcommand{\vv}[1]{\vec{#1}}\newcommand{\ds}{\displaystyle}
\newcommand{\R}{\mathbb{R}}\newcommand{\N}{\mathbb{N}}\newcommand{\Z}{\mathbb{Z}}
% (un \usepackage{fancyhdr}+en-tête est possible ; il est ignoré côté web)
% =========================================================================
\begin{document}
\sisetup{per-mode=symbol}

\begin{corrige}[variations composées de la force]
\begin{enumerate}[label=\alph*)]
  \item Doubler chaque charge multiplie $F$ par $2\times 2 = 4$ ; diviser $d$ par $2$ multiplie $F$
        par $2^2 = 4$. Au total $\times 16$ :
        \[ F = 16\,F_0 = \SI{16}{\newton}. \]
  \item Diviser chaque charge par $3$ : facteur $\tfrac{1}{9}$ ; diviser $d$ par $3$ : facteur
        $3^2 = 9$. Les deux se compensent exactement :
        \[ F = \frac{1/9}{1/9}\,F_0 = F_0 = \SI{1}{\newton}\ \text{(inchangée)}. \]
\end{enumerate}
\end{corrige}

\begin{corrige}[deux forces perpendiculaires]
\begin{enumerate}[label=\alph*)]
  \item Pythagore : $F = \sqrt{F_1^2 + F_2^2} = \sqrt{6^2 + 8^2} = \sqrt{36+64} = \sqrt{100} = \SI{10}{\newton}$.
  \item $\tan\theta = \dfrac{F_2}{F_1} = \dfrac{8}{6} = \dfrac{4}{3}$, d'où $\theta = \arctan(4/3) \approx \ang{53.1}$.
\end{enumerate}
\end{corrige}

\begin{corrige}[triangle équilatéral]
\begin{enumerate}[label=\alph*)]
  \item Les deux forces ont même intensité $f = K q^2/a^2$ et font entre elles un angle de
        $\ang{60}$. Par symétrie la résultante est verticale, somme des projections sur la
        bissectrice :
        \[ F_{\text{tot}} = 2 f\cos\ang{30} = 2\,\frac{K q^2}{a^2}\times\frac{\sqrt3}{2}
             = \frac{\sqrt3\,K q^2}{a^2}. \]
  \item $f = \num{9e9}\times\dfrac{(\num{2e-6})^2}{(0{,}10)^2} = \num{9e9}\times\dfrac{\num{4e-12}}{0{,}01} = \SI{3.6}{\newton}$, donc
        \[ F_{\text{tot}} = \sqrt3\times 3{,}6 \approx \SI{6.2}{\newton}. \]
\end{enumerate}
\end{corrige}

\begin{corrige}[trois charges alignées]
\begin{enumerate}[label=\alph*)]
  \item Les distances $A\!B$ et $B\!C$ valent $\SI{0.10}{\metre}$ et les produits de charges sont
        identiques, donc les deux forces ont la même intensité :
        \[ F_{A/B} = F_{C/B} = \num{9e9}\times\frac{(\num{2e-6})(\num{3e-6})}{(0{,}10)^2}
              = \num{9e9}\times\frac{\num{6e-12}}{0{,}01} = \SI{5.4}{\newton}. \]
        $A$ ($+$) \emph{repousse} $B$ : force vers la \textbf{droite}. $C$ ($-$) \emph{attire} $B$ :
        force aussi vers la \textbf{droite}.
  \item Les deux forces ont le \emph{même sens} (vers la droite) : elles s'ajoutent.
        \[ F_{\text{tot}} = 5{,}4 + 5{,}4 = \SI{10.8}{\newton}\ \text{vers la droite (vers } C). \]
\end{enumerate}
\end{corrige}

\begin{corrige}[superposition à angle droit : calcul complet]
\begin{enumerate}[label=\alph*)]
  \item \[ F_1 = \num{9e9}\times\frac{(\num{2e-6})(\num{2e-6})}{(0{,}20)^2}
             = \num{9e9}\times\frac{\num{4e-12}}{0{,}04} = \SI{0.9}{\newton}\ (\text{horizontale}), \]
        \[ F_2 = \num{9e9}\times\frac{(\num{2e-6})(\num{6e-6})}{(0{,}30)^2}
             = \num{9e9}\times\frac{\num{1.2e-11}}{0{,}09} = \SI{1.2}{\newton}\ (\text{verticale}). \]
  \item Forces perpendiculaires :
        \[ F = \sqrt{F_1^2 + F_2^2} = \sqrt{0{,}9^2 + 1{,}2^2} = \sqrt{0{,}81 + 1{,}44}
             = \sqrt{2{,}25} = \SI{1.5}{\newton}. \]
        (Les deux charges $q_1,q_2$ étant positives, chaque force \emph{repousse} $q_0$ : la
        résultante pointe à l'opposé du coin $(q_1,q_2)$.)
\end{enumerate}
\end{corrige}

\end{document}
